% Keep the horizon evidence immediately beside its discussion.
\begingroup
\setlength{\intextsep}{8pt}
\begin{table}[H]
\caption{\PaperRevision{Effect of the look-ahead range on compilation results}}
\label{tab:sensitivity}
\centering
\PaperTableSetup
\fontsize{8.5pt}{9.7pt}\selectfont
\setlength{\tabcolsep}{2pt}
\begin{tabular*}{\columnwidth}{@{}l@{\extracolsep{\fill}}rrrrr@{}}
\toprule
Look-ahead upper bound $H$ & 0 & 1 & 2 & 4 & 8\\
\midrule
\PaperRevision{Fidelity ratio} & \HorizonExtZeroFidelityRatio & \HorizonExtOneFidelityRatio & \HorizonExtTwoFidelityRatio & \HorizonExtFourFidelityRatio & \HorizonExtEightFidelityRatio\\
\PaperRevision{Batch ratio} & \HorizonExtZeroBatchRatio & \HorizonExtOneBatchRatio & \HorizonExtTwoBatchRatio & \HorizonExtFourBatchRatio & \HorizonExtEightBatchRatio\\
\bottomrule
\end{tabular*}
\PaperTableNote{\fontsize{8pt}{9.1pt}\selectfont\PaperRevision{All ratios use the same circuit's $H=8$ result as a reference of 1; $H=0$ evaluates only the current layer. Of \HorizonExtTargetN{} circuits, \HorizonExtCompleteN{} meet the comparison criteria under all five configurations. For each circuit and configuration, we take the median of \HorizonExtSeedN{} runs, divide by the circuit's $H=8$ median, and then take the geometric mean. Values are rounded to four decimal places.}}
\end{table}
\endgroup
